\section{Introduction}

The advent of large-scale quantum computing poses a fundamental threat to the
public-key cryptography underpinning today's digital infrastructure.
Shor's algorithm~\cite{SIAM:Shor97} demonstrates that a sufficiently powerful
quantum computer can factor large integers and compute discrete logarithms in
polynomial time, thereby rendering RSA and elliptic-curve cryptography (ECC)
insecure.
Grover's algorithm~\cite{STOC:Grover96} further reduces the effective security
of symmetric ciphers and hash functions by providing a quadratic speedup over
classical brute-force search, effectively halving key lengths.
While the timeline to cryptographically relevant quantum computers remains
uncertain, the strategic risk is real: adversaries may harvest ciphertext today
and decrypt it once quantum capability matures—a threat model known as
``harvest now, decrypt later.''
In response, the National Institute of Standards and Technology (NIST) has
finalized two post-quantum cryptography (PQC) standards: FIPS~203 (ML-KEM,
based on CRYSTALS-Kyber)~\cite{NIST:FIPS203} for key encapsulation, and
FIPS~204 (ML-DSA, based on CRYSTALS-Dilithium)~\cite{NIST:FIPS204} for digital
signatures.
NIST IR~8547~\cite{NIST:IR8547} further mandates that federal agencies complete
migration away from RSA and ECC by 2035.
Physical access control systems are a particularly sensitive target: NFC-based
door locks and vehicle entry systems guard critical infrastructure and personal
safety, yet virtually all deployed protocols rely on ECC.
Transitioning these systems to post-quantum security before the quantum threat
materialises is therefore an urgent engineering challenge.

Aliro~1.0~\cite{CSA:Aliro10}, published by the Connectivity Standards Alliance
(CSA), is a recently standardised multi-modal access-control protocol supporting
NFC, BLE, and UWB bearers.
Its cryptographic core is built on ECC: an ECDH key-exchange over P-256 for
session-key establishment and ECDSA for certificate-based entity authentication.
Similar designs appear in the automotive domain through CCC Digital
Key~3.0~\cite{CCC:DigKey3}, which likewise relies on ECDH/ECDSA for secure
vehicle-phone pairing.
Integrating ML-KEM and ML-DSA into such protocols is non-trivial for three
interconnected reasons.
First, \emph{object-size inflation}: an ML-DSA-65 public key occupies 1{,}952
bytes and a signature 3{,}309 bytes, compared with 65 bytes and $\approx$72
bytes for ECDSA on P-256---an increase of roughly $30\times$ and $46\times$,
respectively.
An ML-KEM-768 ciphertext is 1{,}088 bytes, versus 65 bytes for an ephemeral
ECDH public key.
Second, \emph{NFC channel constraints}: ISO/IEC~14443 ISO-DEP frames carry at
most 255 bytes of application data per APDU, and the physical-layer bit rate
ranges from 106 to 848~kbit/s.
Even at 424~kbit/s---the rate typically negotiated by Android HCE---a single
PQ-Aliro transaction under the default ML-DSA-65 $\times$ ML-KEM-768
configuration transfers approximately 16.4~KB across $\approx$75 APDUs, roughly
$14\times$ the $\approx$1.2~KB of a classical ECC transaction.
Third, \emph{platform limitations}: Android's \texttt{KeyStore} API exposes no
native support for lattice-based algorithms, requiring the implementation to
manage PQC key material entirely in software using third-party libraries such as
BouncyCastle.

This paper presents \textbf{PQ-Aliro}, the first complete integration of
NIST-standardised post-quantum algorithms into the Aliro NFC access-control
protocol.
Our contributions are as follows:

\begin{enumerate}

\item \textbf{Protocol design and implementation.}
We design a drop-in PQC variant of the Aliro Expedited flow that replaces
ECDH with ML-KEM-768 key encapsulation and ECDSA with ML-DSA-65 signatures,
preserving the three-round AUTH0/LOADCERT/AUTH1 message structure.
We provide a full open-source implementation: an Android HCE user-device app
and a PC NFC reader, both using BouncyCastle~1.78 BCPQC as the cryptographic
back-end.

\item \textbf{APDU-level overhead analysis.}
We conduct an exhaustive byte-accurate analysis of communication overhead for
all nine PQC parameter combinations (Dilithium$\{2,3,5\} \times
$Kyber$\{512,768,1024\}$).
The default ML-DSA-65 $\times$ ML-KEM-768 configuration incurs $\approx$16.4~KB
and $\approx$75 APDUs per transaction; the lightest ML-DSA-44 $\times$
ML-KEM-512 configuration reduces this to $\approx$11.9~KB and $\approx$55
APDUs.
We additionally show that a \emph{key-slot} optimisation---pre-caching the
reader certificate on the user device---eliminates the LOADCERT round and saves
$\approx$5--7~KB.

\item \textbf{Formal security analysis.}
Building on the Bellare--Rogaway security model~\cite{C:BelRog93} and the
provable security of ML-KEM~\cite{EuroSP:BDKLLSSSS18} and
ML-DSA~\cite{TCHES:DKLLSSS18}, we provide a formal security proof establishing
that PQ-Aliro achieves mutual authentication, forward secrecy, and
reader-identity privacy under the module learning-with-errors (MLWE) and
module short-integer-solution (MSIS) assumptions (see §\,6--7 for the full model
and proof).

\item \textbf{Engineering insights for constrained Android HCE.}
We document critical implementation findings: isolating Dilithium private-key
operations inside an \texttt{isolatedProcess} Service to prevent cross-process
key leakage, and dispatching ML-DSA signing asynchronously to a background
coroutine to avoid ANR deadlocks on the HCE main thread.
These lessons generalise to any lattice-based protocol deployed over Android HCE.

\end{enumerate}

\noindent
The remainder of this paper is organised as follows.
Section~2 surveys related work on PQC integration in constrained communication
protocols and NFC/BLE access control.
Section~3 provides background on the Aliro~1.0 protocol and the ML-KEM/ML-DSA
algorithms.
Section~4 introduces formal definitions and notation used throughout the security
analysis.
Section~5 presents the PQ-Aliro protocol design, APDU message formats, and trust
framework.
Section~6 defines the security model and security goals.
Section~7 contains the formal security proof.
Section~8 describes the Android HCE and PC reader implementation details.
Section~9 reports the communication-overhead evaluation and NFC timing
measurements.
Section~10 discusses limitations, optimisation trade-offs, and future directions.
Section~11 concludes the paper.
